\SourceLeaf{459}
\noindent Dividing these two equations by one another and then taking the square root gives
\[
\frac{\sin\tfrac12 a^3}{\cos\tfrac12 a}.\frac{\cos\tfrac12 b}{\sin\tfrac12 b^3}
=\frac{\sin(A+\omega)\sin(A-\tfrac12\omega)^2}{\sin(B+\omega)\sin(B-\tfrac12\omega)^2}.
\]
\noindent This equation can also be put in the form
\[
\frac ab=\frac{\sin A}{\sin B}.\mkern6mu{}^3\mkern-4mu\BareRoot{D},
\]
\Needspace{10\baselineskip}
\noindent where $D$ is written in place of
\[
\begin{aligned}
&\frac{a^3\cos\tfrac12 a}{8\sin\tfrac12 a^3}.\frac{8\sin\tfrac12 b^3}{b^3\cos\tfrac12 b}\, .\\[5pt]
&\quad\frac{\sin(A+\omega)\sin(A-\tfrac12\omega)^2}{\sin A^3}.\frac{\sin B^3}{\sin(B+\omega)\sin(B-\tfrac12\omega)^2}
\end{aligned}
\]
\noindent for brevity.

This formula is rigorously exact. But it is immediately apparent that if $a,b,c$ are very small and are regarded as quantities of the first order, each of the four factors of which $D$ is composed differs from unity only by quantities of the fourth order.

This subject is treated according to more general principles and extended to triangles formed by shortest lines on arbitrary curved surfaces in my \textit{Disquisitiones generales circa superficies curvas.}

